\section{Regionen und Verfügbarkeitszonen}
\subsection{Region}
Eine Cloud-Region ist ein geografischer Bereich, der aus mehreren
Verfügbarkeitszonen besteht. Regionen sind meist nach Himmelsrichtung
oder Lage benannt, zum Beispiel \textit{Deutschland West Central} oder
\textit{North Europe} bei Azure. Die Wahl der Region bestimmt, wo Daten
physisch gespeichert werden (Datenresidenz) und beeinflusst die Latenz
zu den Endnutzern. Jede Region ist physisch und logisch von anderen
Regionen isoliert, sodass ein Ausfall einer Region andere nicht
beeinträchtigt.
\subsection{Verfügbarkeitszonen}
Eine Verfügbarkeitszone (engl. \textit{Availability Zone}) ist eine
physisch getrennte Einheit innerhalb einer Region. Sie besteht aus einem
oder mehreren Rechenzentren mit eigener Stromversorgung, Kühlung und
Netzwerkanbindung. Verfügbarkeitszonen innerhalb einer Region sind über
redundante Hochgeschwindigkeitsverbindungen mit geringer Latenz
verbunden. Durch die Verteilung von Diensten auf mehrere Zonen können
Anwendungen auch bei einem Zonenausfall verfügbar bleiben
(Hochverfügbarkeit).
